\section{Identifiers Usage}
\label{sec:grust:id_usage}

This section formally defines the rules governing the use of identifiers in \GRust. A precise and
unambiguous specification of identifier scoping is essential for any programming language, as it
dictates where an identifier can be used and what it refers to. For a language like \GRust, intended
for critical systems, these rules are part of the language definition, ensuring predictable and
verifiable behavior.

The formal system presented here serves as a blueprint for the static analysis performed by the
compiler. It provides a rigorous foundation for implementing the symbol table, which tracks all
identifiers and their associated information, such as scope and type. By adhering to these rules,
the compiler can detect a range of errors at compile-time, including the use of undefined
identifiers, the redeclaration of an identifier within the same scope, and more subtle violations
specific to the structure of \GRust. For example, the logic enforces that all branches of a
\kft{match} equation must define the same set of identifiers, guaranteeing accessible values for
every execution.

This formalization addresses the two distinct scoping models within \GRust: the sequential, lexical
scoping of \GRFRP services and functions, and the mutually recursive scoping of \GRSync components.
The following logic, expressed through a set of inference rules, details how identifiers are
declared, defined, and used correctly across all language constructs.

\begin{table}[!ht]
    \centering
    \begin{tabularx}{\textwidth}{|l|X|}\hline
        $\ID_s, \: \ID_f, \: \ID_c, \: \ID$                                             & Service, function,
        component, and variable identifiers respectively.                                                                        \\\hline
        $\IDs ::= \{\ID^*\}$                                                            & A set of identifiers.                  \\\hline
        \begin{minipage}{0.37\textwidth}
            $\scope ::= \importScope \: \vert \: \exportScope \:
                \vert \:\inputScope$ \\ $\vert \: \outputScope \: \vert
                \: \localScope \: \vert \: \memoryScope$
        \end{minipage}                         & The scope of an identifier, such as an import, export,
        input, output, local, or component memory.                                                                               \\\hline
        $\G[\ID_s]$, $\G[\ID_f]$, $\G[\ID_c]$                                           & The identifier context maps
        service, function, or component identifiers to their definitions.                                                        \\\hline
        $\G[\ID] = \scope$                                                              & The identifier context maps a variable
        identifier to its scope.                                                                                                 \\\hline
        \idUsage{\G}{elem}{\texttt{Elem}}                                               & Asserts the correct usage
        of identifiers in the \GRust construct \texttt{Elem} with respect to the context \G.
        \\\hline
        \exports{(\Prog \: \vert \: \Item)}{\IDs}                                       & \IDs is the set of
        identifiers exported by program \Prog or item \Item.                                                                     \\\hline
        \declares{(\FlowStmt \: \vert \: \Eq \: \vert \: \Pat \: \vert \: \EPat)}{\IDs} & \IDs is the set of
        identifiers declared by statement \FlowStmt, equation \Eq or patterns \Pat and \EPat.                                    \\\hline
        \inits{\Eq}{\IDs}                                                               & \IDs is the set of
        identifiers initialized by equation \Eq.                                                                                 \\\hline
        \defines{(\Prog \: \vert \: \Item \: \vert \: \FlowStmt
        \: \vert \: \Eq)}{\IDs}                                                         & \IDs is the set of
        identifiers defined by program \Prog, item \Item, statement \FlowStmt, or equation \Eq, including declared ones.         \\\hline
    \end{tabularx}
    \caption[Notations for the usage of identifiers in \GRust.]%
    {Notations for the logic of identifier usage in \GRust.}
    \label{tab:grust:id_usage:notations}
\end{table}

We first introduce the notations for this logic, summarized in \Cref{tab:grust:id_usage:notations}.
The core of our formal system is the judgment \idUsage{\G}{elem}{\texttt{Elem}}, which asserts that
a \GRust construct \texttt{Elem} uses identifiers correctly with respect to an identifier context
\G. The context \G acts as a symbol table, mapping identifiers to their \emph{scope}, such as
\importScope, \outputScope, or \localScope. It also maps service, function, and component
identifiers ($\ID_s, \ID_f, \ID_c$) to their definitions.

To support this central judgment, we define several helper judgments that extract sets of
identifiers (\IDs) from various constructs:
\begin{itemize}
    \item \exports{\Prog \: \vert \: \Item}{\IDs} retrieves the set of identifiers that a program or
          item exports to the external environment.
    \item \declares{(\FlowStmt \: \vert \: \Eq \: \vert \: \Pat \: \vert \: \EPat)}{\IDs} retrieves
          the set of identifiers newly declared within a statement, an equation, or pattern.
    \item \inits{\Eq}{\IDs} retrieves the set of memory identifiers initialized in an equation.
    \item \defines{(\Prog \: \vert \: \Item \: \vert \: \FlowStmt \: \vert \: \Eq)}{\IDs} retrieves
          the set of identifiers that are given a definition. This includes both newly declared
          identifiers and existing ones (\eg component outputs).
\end{itemize}
The following paragraphs detail the rules for these judgments, explaining the logic of identifier
usage for each part of the \GRust language.

\subsection{\GRust Programs}
\label{sec:grust:id_usage:prog}

A \GRust program is well-formed if all its items are well-formed. The rule
\nameref{rule:grust:id_usage:program} formalizes this by checking each item individually. Crucially,
it also enforces a key property of a complete program: the set of identifiers it \emph{exports} must
be identical to the set it \emph{defines}. This ensures that every declared export has a
corresponding definition within a service, and that there are no definitions for non-exported flows
at the top level.

\begin{center}
    \begin{minipage}{0.9\textwidth}
        \centering
        \begin{mathpar}
            %% Program
            \inferID
            {
                \left(\itemIdUsage{\Item}\right)^* \\
                \exports{\Item^*}{\IDs} \\
                \defines{\Item^*}{\IDs}
            }
            { \progIdUsage{\Item^*} }
            {Prog}
            \label{rule:grust:id_usage:program}
        \end{mathpar}
    \end{minipage}
\end{center}

The good usage of an identifier in an item \Item with respect to the identifier context \G is
denoted by \itemIdUsage{\Item}. For interfaces, rules
\nameref{rule:grust:id_usage:interface:import} and
\nameref{rule:grust:id_usage:interface:export} ensure that imported and exported identifiers are
present in the context with the correct scope (\importScope for imports, and \exportScope for
exports).

\begin{center}
    \begin{minipage}{0.9\textwidth}
        \centering
        \begin{mathpar}
            %% Interface
            % Import
            \inferID
            { \G[\ID] = \importScope }
            { \itemIdUsage{\importX{\typed{\FlowKind}{\ID}{\TY}}} }
            {Import}
            \label{rule:grust:id_usage:interface:import}

            % Export
            \inferID
            { \G[\ID] = \exportScope }
            { \itemIdUsage{\exportX{\typed{\FlowKind}{\ID}{\TY}}} }
            {Export}
            \label{rule:grust:id_usage:interface:export}
        \end{mathpar}
    \end{minipage}
\end{center}

Rule \nameref{rule:grust:id_usage:service} states that for a service to be correct, its definition
must be in the context and its body must adhere to the identifier usage rules.

\begin{center}
    \begin{minipage}{0.9\textwidth}
        \centering
        \begin{mathpar}
            %% Service
            \inferID
            {
                \G[\ID_s] = \servItem{\ID_s}{\interval{\INT}{\INT}}{\many{\FlowStmt}} \\\\
                \stmtIdUsage{\G}{\many{\FlowStmt}}
            }
            { \itemIdUsage{\servItem{\ID_s}{\interval{\INT}{\INT}}{\many{\FlowStmt}}} }
            {Serv}
            \label{rule:grust:id_usage:service}
        \end{mathpar}
    \end{minipage}
\end{center}

For functions and components, the logic is similar. The rules ensure the item definition exists in
the context and that the body is well-formed. For a function, applying rule
\nameref{rule:grust:id_usage:fun}, the context is extended with its inputs before checking the body.
For a component, applying rule \nameref{rule:grust:id_usage:comp}, the context is extended with both
inputs and outputs to support the mutually recursive nature of its body.

\begin{center}
    \begin{minipage}{0.9\textwidth}
        \centering
        \begin{mathpar}
            %% Function
            \inferID
            {
                \G[\ID_f] = \funItem{\ID_f}{\many{(\ID_i \: \kft{:} \: \TY_i)}}{\TY_o}{\many{\Decl}} \\\\
                \G' = \G \ctxconcat [(\inputScope / \ID_i)^*] \\
                \declIdUsage{\G'}{\many{\Decl} \: \returnStmt{\Expr}}
            }
            { \itemIdUsage{\funItem{\ID_f}{\many{(\ID_i \: \kft{:} \: \TY_i)}}{\TY_o}{\many{\Decl}}} }
            {Fun}
            \label{rule:grust:id_usage:fun}
        \end{mathpar}
    \end{minipage}
\end{center}


\begin{center}
    \begin{minipage}{0.9\textwidth}
        \centering
        \begin{mathpar}
            %% Component
            \inferID
            {
                \G[\ID_c] = \compItem{\ID_c}{\many{(\ID_i \: \kft{:} \: \STY_i)}}{\many{(\ID_o \: \kft{:} \: \STY_o)}}{\many{\Eq}} \\\\
                \G' = \G \ctxconcat [(\inputScope / \ID_i)^*] \ctxconcat [(\outputScope / \ID_o)^*] \\\\
                \eqIdUsage{\G'}{\setEq} \\
                \defines{\setEq}{\IDs} \\
                \left(\ID_o \in \IDs\right)^*
            }
            { \itemIdUsage{\compItem{\ID_c}{\many{(\ID_i \: \kft{:} \: \STY_i)}}{\many{(\ID_o \: \kft{:} \: \STY_o)}}{\many{\Eq}}} }
            {Comp}
            \label{rule:grust:id_usage:comp}
        \end{mathpar}
    \end{minipage}
\end{center}

\paragraph{Exported identifiers}
\label{sec:grust:id_usage:prog:exported}

The judgment \exports{\Prog}{\IDs} computes the set of identifiers a program makes available
externally. As shown in rule \nameref{rule:grust:id_usage:prog:exported:program}, this is the union
of identifiers from all \kft{export} items. The use of the \emph{disjoint union} operator
($\bigsqcup$) enforces that each exported identifier must have a unique name across the entire
program, preventing ambiguity at the system level. The subsequent rules show that only \kft{export}
items contribute to this set.

\begin{center}
    \begin{minipage}{0.9\textwidth}
        \centering
        \begin{mathpar}
            %% Program
            \inferEXP
            { \left(\exports{\Item}{\IDs}\right)^* }
            { \exports{\Item^*}{\bigsqcup \IDs^*} }
            {Prog}
            \label{rule:grust:id_usage:prog:exported:program}
        \end{mathpar}
    \end{minipage}

    \begin{minipage}{0.99\textwidth}
        \centering
        \begin{mathpar}
            %% Interface
            % Import
            \inferEXP
            { }
            { \exports{\importX{\typed{\FlowKind}{\ID}{\TY}}}{\emptyset} }
            {Import}
            \label{rule:grust:id_usage:prog:exported:interface:import}

            % Export
            \inferEXP
            { }
            { \exports{\exportX{\typed{\FlowKind}{\ID}{\TY}}}{\{\ID\}} }
            {Export}
            \label{rule:grust:id_usage:prog:exported:interface:export}

            %% Service
            \inferEXP
            { }
            { \exports{\servItem{\ID_s}{\interval{\INT}{\INT}}{\many{\FlowStmt}}}{\emptyset} }
            {Serv}
            \label{rule:grust:id_usage:prog:exported:service}

            %% Function
            \inferEXP
            { }
            { \exports{\funItem{\ID_f}{\many{(\ID_i \: \kft{:} \: \TY_i)}}{\TY_o}{\many{\Decl}}}{\emptyset} }
            {Fun}
            \label{rule:grust:id_usage:prog:exported:fun}

            %% Component
            \inferEXP
            { }
            { \exports{\compItem{\ID_c}{\many{(\ID_i \: \kft{:} \: \STY_i)}}{\many{(\ID_o \: \kft{:} \: \STY_o)}}{\many{\Eq}}}{\emptyset} }
            {Comp}
            \label{rule:grust:id_usage:prog:exported:comp}
        \end{mathpar}
    \end{minipage}
\end{center}

\paragraph{Defined identifiers}
\label{sec:grust:id_usage:prog:defined}

The judgment \defines{\Prog}{\IDs} computes the set of identifiers defined within the services of a
program. Rule \nameref{rule:grust:id_usage:prog:defined:program} aggregates these identifiers from
all items, again using a disjoint union ($\bigsqcup$) to ensure that each defined identifier is
unique. As specified by the rules below, only instantiation statements (for exported flows) inside a
service contribute to this set.

\begin{center}
    \begin{minipage}{0.9\textwidth}
        \centering
        \begin{mathpar}
            \inferDEF
            { \left(\defines{\Item}{\IDs}\right)^* }
            { \defines{\Item^*}{\bigsqcup \IDs^*} }
            {Prog}
            \label{rule:grust:id_usage:prog:defined:program}
        \end{mathpar}
    \end{minipage}

    \begin{minipage}{0.9\textwidth}
        \centering
        \begin{mathpar}
            %% Interface
            % Import
            \inferDEF
            { }
            { \defines{\importX{\typed{\FlowKind}{\ID}{\TY}}}{\emptyset} }
            {Import}
            \label{rule:grust:id_usage:prog:defined:interface:import}

            % Export
            \inferDEF
            { }
            { \defines{\exportX{\typed{\FlowKind}{\ID}{\TY}}}{\emptyset} }
            {Export}
            \label{rule:grust:id_usage:prog:defined:interface:export}
        \end{mathpar}
    \end{minipage}

    \begin{minipage}{0.9\textwidth}
        \centering
        \begin{mathpar}
            %% Service
            \inferDEF
            {
                \declares{\block{\many{\FlowStmt}}}{\IDs_{\texttt{let}}} \\
                \defines{\block{\many{\FlowStmt}}}{\IDs_{\texttt{let}} \sqcup \IDs_{\texttt{export}}}
            }
            { \defines{\servItem{\ID_s}{\interval{\INT}{\INT}}{\many{\FlowStmt}}}{\IDs_{\texttt{export}}} }
            {Serv}
            \label{rule:grust:id_usage:prog:defined:service}

        \end{mathpar}
    \end{minipage}

    \begin{minipage}{0.9\textwidth}
        \centering
        \begin{mathpar}
            %% Function
            \inferDEF
            { }
            { \defines{\funItem{\ID_f}{\many{(\ID_i \: \kft{:} \: \TY_i)}}{\TY_o}{\many{\Decl}}}{\emptyset} }
            {Fun}
            \label{rule:grust:id_usage:prog:defined:fun}

        \end{mathpar}
    \end{minipage}

    \begin{minipage}{0.9\textwidth}
        \centering
        \begin{mathpar}
            %% Component
            \inferDEF
            { }
            { \defines{\compItem{\ID_c}{\many{(\ID_i \: \kft{:} \: \STY_i)}}{\many{(\ID_o \: \kft{:} \: \STY_o)}}{\many{\Eq}}}{\emptyset} }
            {Comp}
            \label{rule:grust:id_usage:prog:defined:comp}
        \end{mathpar}
    \end{minipage}
\end{center}

\subsection{Bodies of \GRFRP Services}
\label{sec:grust:id_usage:grfrp}

This part describes the correct usage of identifiers in services and all constructs of \GRFRP.

\subsubsection{Flow Statements}
\label{sec:grust:id_usage:grfrp:flow_stmt}

Flow statements in \GRFRP services follow a sequential, lexical scoping discipline. Identifiers are
introduced and are only visible to the statements that follow them in the same block. This is
captured by the judgment \stmtIdUsage{\G}{\FlowStmt^*}, which threads the identifier context \G
through the list of statements.

An empty block of statements is trivially well-formed, as stated by rule
\nameref{rule:grust:id_usage:grfrp:flow_stmt:skip}.
A \kft{let} statement introduces new local identifiers: rule
\nameref{rule:grust:id_usage:grfrp:flow_stmt:let} checks the right-hand side of the declaration
before extending the context \G with these new identifiers for the subsequent statements. This
ensures the new identifiers are available in the rest of the scope but not in their own definition.
In contrast, a definition statement assigns a flow to an existing, exported identifier: rule
\nameref{rule:grust:id_usage:grfrp:flow_stmt:def} verifies that the identifier has already been
declared with \exportScope and then checks the remaining statements with the original context.

\begin{center}
    \begin{minipage}{0.9\textwidth}
        \centering
        \begin{mathpar}
            %% Flow Statement
            % Skip
            \inferIDFlowStmt
            { }
            { \stmtIdUsage{\G}{} }
            {Skip}
            \label{rule:grust:id_usage:grfrp:flow_stmt:skip}

            % Let
            \inferIDFlowStmt
            {
                \flowIdUsage{\G}{\FlowOp} \\
                \G' = \G \ctxconcat [(\localScope / \ID)^*] \\
                \stmtIdUsage{\G'}{\FlowStmt^*}
            }
            { \stmtIdUsage{\G}{\letStmt{\many{(\typed{\FlowKind}{\ID}{\TY})}}{\FlowOp} \: \FlowStmt^*} }
            {Let}
            \label{rule:grust:id_usage:grfrp:flow_stmt:let}

            % Definition
            \inferIDFlowStmt
            {
                \left(\G[\ID] = \exportScope\right)^* \\
                \flowIdUsage{\G}{\FlowOp} \\
                \stmtIdUsage{\G}{\FlowStmt^*}
            }
            { \stmtIdUsage{\G}{\defStmt{\many{\ID}}{\FlowOp} \: \FlowStmt^*} }
            {Def}
            \label{rule:grust:id_usage:grfrp:flow_stmt:def}
        \end{mathpar}
    \end{minipage}
\end{center}

\paragraph{Declared identifiers}
\label{sec:grust:id_usage:grfrp:flow_stmt:declared}

The set of declared identifiers in a flow statement is computed by the judgment
\declares{\FlowStmt}{\IDs}. Rules \nameref{sec:grust:id_usage:grfrp:flow_stmt:declared:let}
and \nameref{sec:grust:id_usage:grfrp:flow_stmt:declared:def} ensure that only \kft{let}
statements contribute to the set of declared identifiers at program-level.

\begin{center}
    \begin{minipage}{0.9\textwidth}
        \centering
        \begin{mathpar}
            \inferDECLFlowStmt
            { }
            { \declares{\letStmt{\many{(\ldots)}}{\ldots} \: \FlowStmt^*}{\{\many{\ID}\}} }
            {Let}
            \label{sec:grust:id_usage:grfrp:flow_stmt:declared:let}

            \inferDECLFlowStmt
            { }
            { \declares{\defStmt{\many{\ID}}{\ldots} \: \FlowStmt^*}{\emptyset} }
            {Def}
            \label{sec:grust:id_usage:grfrp:flow_stmt:declared:def}
        \end{mathpar}
    \end{minipage}
\end{center}

\paragraph{Defined identifiers}
\label{sec:grust:id_usage:grfrp:flow_stmt:defined}

The set of defined identifiers in a flow statement is computed by the judgment
\defines{\FlowStmt}{\IDs}. Rules \nameref{sec:grust:id_usage:grfrp:flow_stmt:defined:let}
and \nameref{sec:grust:id_usage:grfrp:flow_stmt:defined:def} specify that all statements contribute
to the set of defined identifiers for the program-level check.

\begin{center}
    \begin{minipage}{0.9\textwidth}
        \centering
        \begin{mathpar}
            \inferDEFFlowStmt
            { }
            { \defines{\letStmt{\many{(\typed{\FlowKind}{\ID}{\TY})}}{\FlowOp} \: \FlowStmt^*}{\{\many{\ID}\}} }
            {Let}
            \label{sec:grust:id_usage:grfrp:flow_stmt:defined:let}

            \inferDEFFlowStmt
            { }
            { \defines{\defStmt{\many{\ID}}{\FlowOp} \: \FlowStmt^*}{\{\many{\ID}\}} }
            {Def}
            \label{sec:grust:id_usage:grfrp:flow_stmt:defined:def}
        \end{mathpar}
    \end{minipage}
\end{center}

\subsubsection{Flow Operations}
\label{sec:grust:id_usage:grfrp:flow_op}

The judgment \flowIdUsage{\G}{\FlowOp} checks that a flow operation uses identifiers
correctly. The rules ensure that any used identifier, function, or component is present in
the context \G and that all sub-expressions are themselves well-formed.

\begin{center}
    \begin{minipage}{0.99\textwidth}
        \centering
        \begin{mathpar}
            %% Flow Operation
            % Identifier
            \inferIDFlow
            { \G[\ID] = \scope }
            { \flowIdUsage{\G}{\ID} }
            {ID}
            \label{rule:grust:id_usage:grfrp:flow_op:ident}

            %% Application
            % Function
            \inferIDFlow
            {
                \G[\ID_f] = \funItem{\ID_f}{\many{(\ID_i \: \kft{:} \: \TY_i)}}{\TY_o}{\many{\Decl}} \\\\
                \left(\flowIdUsage{\G}{\FlowOp}\right)^*
            }
            { \flowIdUsage{\G}{\compApp{\ID_f}{\many{\FlowOp}}} }
            {AppF}
            \label{rule:grust:id_usage:grfrp:flow_op:funapp}

            % Component
            \inferIDFlow
            {
                \G[\ID_c] = \compItem{\ID_c}{\many{(\ID_i \: \kft{:} \: \STY_i)}}{\many{(\ID_o \: \kft{:} \: \STY_o)}}{\many{\Eq}} \\\\
                \left(\flowIdUsage{\G}{\FlowOp}\right)^*
            }
            { \flowIdUsage{\G}{\compApp{\ID_c}{\many{\FlowOp}}} }
            {AppC}
            \label{rule:grust:id_usage:grfrp:flow_op:compapp}
        \end{mathpar}
    \end{minipage}
\end{center}

\GReact operations are checked by ensuring their input flow operations are well-formed.

\begin{center}
    \begin{minipage}{0.9\textwidth}
        \centering
        \begin{mathpar}
            %% GReact
            % OnChange
            \inferIDFlow
            { \flowIdUsage{\G}{\FlowOp} }
            { \flowIdUsage{\G}{\onchangeExpr{\FlowOp}} }
            {OChg}
            \label{rule:grust:id_usage:grfrp:flow_op:onchange}

            % Persist
            \inferIDFlow
            { \flowIdUsage{\G}{\FlowOp} }
            { \flowIdUsage{\G}{\persistExpr{\FlowOp}} }
            {Prst}
            \label{rule:grust:id_usage:grfrp:flow_op:persist}

            % Scan
            \inferIDFlow
            { \flowIdUsage{\G}{\FlowOp} }
            { \flowIdUsage{\G}{\scanExpr{\FlowOp}{\INT}} }
            {Scan}
            \label{rule:grust:id_usage:grfrp:flow_op:scan}

            % Sample
            \inferIDFlow
            { \flowIdUsage{\G}{\FlowOp} }
            { \flowIdUsage{\G}{\scanExpr{\FlowOp}{\INT}} }
            {Samp}
            \label{rule:grust:id_usage:grfrp:flow_op:sample}

            % Timeout
            \inferIDFlow
            { \flowIdUsage{\G}{\FlowOp} }
            { \flowIdUsage{\G}{\timeoutExpr{\FlowOp}{\INT}} }
            {TOut}
            \label{rule:grust:id_usage:grfrp:flow_op:timeout}

            % Merge
            \inferIDFlow
            {
                \flowIdUsage{\G}{\FlowOp_1} \\
                \flowIdUsage{\G}{\FlowOp_2}
            }
            { \flowIdUsage{\G}{\mergeExpr{\FlowOp_1}{\FlowOp_2}} }
            {Mrg}
            \label{rule:grust:id_usage:grfrp:flow_op:merge}

            % Throttle
            \inferIDFlow
            { \flowIdUsage{\G}{\FlowOp} }
            { \flowIdUsage{\G}{\throttleExpr{\FlowOp}{\DELTA}} }
            {Thrt}
            \label{rule:grust:id_usage:grfrp:flow_op:throttle}

            % Time
            \inferIDFlow
            { }
            { \flowIdUsage{\G}{\timeExpr} }
            {Time}
            \label{rule:grust:id_usage:grfrp:flow_op:time}
        \end{mathpar}
    \end{minipage}
\end{center}

\subsection{Bodies of \GRSync Functions and Components}
\label{sec:grust:id_usage:grsync}

This part describes the correct usage of identifiers in functions and components.

\subsubsection{Function Declarations}
\label{sec:grust:id_usage:grsync:fun_decl}

Declarations in functions, like statements in services, are ordered and lexically scoped. The
judgment \declIdUsage{\G}{\many{\Decl} \: \returnStmt{\Expr}} checks a block of declarations
followed by a return statement. Rule \nameref{rule:grust:id_usage:grsync:fun_decl:decl} extends the
context with newly declared local identifiers before checking the subsequent declarations and the
final return statement with \nameref{rule:grust:id_usage:grsync:fun_decl:return}.

\begin{center}
    \begin{minipage}{0.9\textwidth}
        \centering
        \begin{mathpar}
            %% Return
            \inferIDDecl
            { \exprIdUsage{\G}{\Expr} }
            { \declIdUsage{\G}{\returnStmt{\Expr}} }
            {Ret}
            \label{rule:grust:id_usage:grsync:fun_decl:return}


            %% Declarations
            \inferIDDecl
            {
                \G' = \G \ctxconcat [(\localScope / \ID)^*] \\
                \declIdUsage{\G'}{\many{\Decl} \: \returnStmt{\Expr}}
            }
            { \declIdUsage{\G}{\declEq{\many{(\ID \: \kft{:} \: \STY)}}{\Expr} \: \many{\Decl} \: \returnStmt{\Expr}} }
            {Decl}
            \label{rule:grust:id_usage:grsync:fun_decl:decl}
        \end{mathpar}
    \end{minipage}
\end{center}

\subsubsection{Component Equations}
\label{sec:grust:id_usage:grsync:eq}

In contrast to the sequential nature of services and function, equations within \GRSync components
follow a mutually recursive scoping model. All identifiers declared locally or as state memories
within a component body are visible throughout that entire body, including in their own definitions.
This allows for defining interdependent dataflow equations without regard to textual order.

Rule \nameref{rule:grust:id_usage:grsync:eqs} formalizes this by first collecting all locally
declared identifiers ($\ID_l^*$) and memory identifiers ($\ID_m^*$) from the set of equations
\setEq. It then extends the context \G with these identifiers to create a new context \G', which is
subsequently used to check each individual equation in the set. The use of disjoint union in the
helper judgments ensures that each local and memory identifier is declared at most once.

\begin{center}
    \begin{minipage}{0.9\textwidth}
        \centering
        \begin{mathpar}
            %% Set of equations
            \inferIDEq
            {
                \declares{\setEq}{\{\ID_l^*\}} \\
                \inits{\setEq}{\{\ID_m^*\}} \\
                \G' = \G \ctxconcat [(\localScope / \ID_l)^*] \ctxconcat [(\memoryScope / \ID_m)^*] \\
                \left(\eqIdUsage{\G'}{\Eq}\right)^*
            }
            { \eqIdUsage{\G}{\setEq} }
            {Blck}
            \label{rule:grust:id_usage:grsync:eqs}
        \end{mathpar}
    \end{minipage}
\end{center}

Individual declaration, definition, and initialization equations are checked to ensure they assign
to identifiers of the correct scope (\localScope, \outputScope, and \memoryScope, respectively), as
described by the rules \nameref{rule:grust:id_usage:grsync:eq:decl},
\nameref{rule:grust:id_usage:grsync:eq:def}, and \nameref{rule:grust:id_usage:grsync:eq:init}.

\begin{center}
    \begin{minipage}{0.9\textwidth}
        \centering
        \begin{mathpar}
            %% Declarations
            \inferIDEq
            {
                \left(\G[\ID] = \localScope\right)^* \\
                \exprIdUsage{\G}{\Expr}
            }
            { \eqIdUsage{\G}{\declEq{\many{(\ID \: \kft{:} \: \STY)}}{\Expr}} }
            {Decl}
            \label{rule:grust:id_usage:grsync:eq:decl}

            %% Definition
            \inferIDEq
            {
                \left(\G[\ID] = \outputScope\right)^* \\
                \exprIdUsage{\G}{\Expr}
            }
            { \eqIdUsage{\G}{\defEq{\many{\ID}}{\Expr}} }
            {Def}
            \label{rule:grust:id_usage:grsync:eq:def}

            %% Initialization
            \inferIDEq
            { \G[\lastExpr{\ID}] = \memoryScope }
            { \eqIdUsage{\G}{\initEq{\ID}{\CST}} }
            {Init}
            \label{rule:grust:id_usage:grsync:eq:init}
        \end{mathpar}
    \end{minipage}
\end{center}

The rule for a \kft{match} equation formalizes how identifiers are scoped within a conditional block
that branches on the value of an expression. First, it asserts that the expression being matched
must be valid in the surrounding context. Then, for each branch of the \kft{match}, the rule
specifies that a new, temporary scope is created. This branch-specific scope contains all the
identifiers from the outer scope, plus any new local identifiers that are bound by that branch's
pattern. The guard condition and the equations within the branch body are then validated using this
temporary scope. This ensures that the logic inside a branch can safely use the variables
deconstructed from the pattern, while also guaranteeing that these variables are properly isolated
and do not leak into sibling branches or the surrounding scope.

\begin{center}
    \begin{minipage}{0.9\textwidth}
        \centering
        \begin{mathpar}
            %% Match
            \inferIDEq
            {
                \exprIdUsage{\G}{\Expr} \\
                \left(\declares{\Pat_k}{\{\ID_k^*\}}\right)^* \\
                \left(\G_k = \G \ctxconcat [(\localScope / \ID_k)^*] \right)^* \\
                \left(\exprIdUsage{\G_k}{\Expr_k}\right)^* \\
                \left(\eqIdUsage{\G_k}{\setEq[k]}\right)^*
            }
            { \eqIdUsage{\G}{\matchX{\Expr}{\many{(\normBr{\gPat{\Pat_k}{\Expr_k}}{\setEq[k]})}}} }
            {Mtch}
            \label{rule:grust:id_usage:grsync:eq:match}
        \end{mathpar}
    \end{minipage}
\end{center}

The rule for a \kft{when} equation defines the scoping for event-driven logic, where a block of
equations is executed in response to an event. The logic proceeds as follows: for each branch, it
first verifies that the event pattern is valid, meaning it refers to a known event. It then
identifies any new local variables that are bound by the pattern, which correspond to the data
payload carried by the event. Similar to a \kft{match} equation, a temporary, branch-specific scope
is created that includes these new payload variables. Finally, the guard condition and the body of
equations for that branch are validated within this temporary scope. This formalizes a safe
mechanism for handling event data, ensuring that the payload of an event is only accessible to the
code block designated to process it.

\begin{center}
    \begin{minipage}{0.99\textwidth}
        \centering
        \begin{mathpar}
            %% When
            \inferIDEq
            {
                \left(\epatIdUsage{\G}{\EPat_k}\right)^* \\
                \left(\declares{\EPat_k}{\{\ID_k^*\}}\right)^* \\\\
                \left(\G_k = \G \ctxconcat [(\localScope / \ID_k)^*] \right)^* \\
                \left(\exprIdUsage{\G_k}{\Expr_k}\right)^* \\
                \left(\eqIdUsage{\G_k}{\setEq[k]}\right)^*
            }
            { \eqIdUsage{\G}{\whenEq{\initBr{\many{(\ID \: \kfm{=} \: \CST\kft{;})}}}{\many{(\normBr{\gPat{\EPat_k}{\Expr_k}}{\setEq[k]})}}} }
            {When}
            \label{rule:grust:id_usage:grsync:eq:when}
        \end{mathpar}
    \end{minipage}
\end{center}

\paragraph{Declared identifiers}
\label{sec:grust:id_usage:grsync:eq:declared}

The judgment \declares{\Eq}{\IDs} collects newly introduced identifiers. The use of disjoint union
($\bigsqcup$) in rule \nameref{rule:grust:id_usage:grsync:eq:declared:eqs} ensures that an
identifier is declared only once in a block of equations.

\begin{center}
    \begin{minipage}{0.9\textwidth}
        \centering
        \begin{mathpar}
            %% Set of equations
            \inferDECLEq
            { \left(\declares{\Eq}{\IDs}\right)^* }
            { \declares{\setEq}{\bigsqcup \IDs^*} }
            {Blck}
            \label{rule:grust:id_usage:grsync:eq:declared:eqs}

            %% Declarations
            \inferDECLEq
            { }
            { \declares{\declEq{\many{(\ID \: \kft{:} \: \STY)}}{\Expr}}{\{\ID^*\}} }
            {Decl}
            \label{rule:grust:id_usage:grsync:eq:declared:decl}

            %% Definition
            \inferDECLEq
            { }
            { \declares{\defEq{\many{\ID}}{\Expr}}{\emptyset} }
            {Def}
            \label{rule:grust:id_usage:grsync:eq:declared:def}

            %% Initialization
            \inferDECLEq
            { }
            { \declares{\initEq{\ID}{\CST}}{\emptyset} }
            {Init}
            \label{rule:grust:id_usage:grsync:eq:declared:init}
        \end{mathpar}
    \end{minipage}
\end{center}

The scoping rules for \kft{match} and \kft{when} equations reflect a fundamental distinction in
their purpose.

A \kft{match} equation defines flows based on pattern matching. As defined by rule
\nameref{rule:grust:id_usage:grsync:eq:declared:match}, all branches must declare an identical set
of identifiers. This ensures that a \kft{match} equation is sound: it always computes a complete and
consistent set of outputs. Regardless of which branch is executed, dependent computations can rely
on the same set of identifiers being available, which prevents the possibility of undefined
behavior.

In contrast, a \kft{when} equation is designed for event-driven \emph{updates}. Rule
\nameref{rule:grust:id_usage:grsync:eq:declared:when} allows each branch to declare a different set
of identifiers, enabling partial state changes. If a signal is not declared in a branch, its value
is not updated when that branch is active and retains its previous value. This persistence model
explains why signals in a \kft{when} construct must be initialized. Similarly, an event is emitted
only if it is declared in the active branch.

\begin{center}
    \begin{minipage}{0.95\textwidth}
        \centering
        \begin{mathpar}
            %% Match
            \inferDECLEq
            { \left(\declares{\setEq[k]}{\IDs}\right)^* }
            { \declares{\matchX{\Expr}{\many{(\normBr{\gPat{\Pat_k}{\Expr_k}}{\setEq[k]})}}}{\IDs} }
            {Mtch}
            \label{rule:grust:id_usage:grsync:eq:declared:match}

            %% When
            \inferDECLEq
            { \left(\declares{\setEq[k]}{\IDs}\right)^* }
            { \declares{
                    \whenEq{\initBr{\many{(\ID \: \kfm{=} \: \CST\kft{;})}}}
                    {\ldots}
                }{\bigcup \IDs^*} }
            {When}
            \label{rule:grust:id_usage:grsync:eq:declared:when}
        \end{mathpar}
    \end{minipage}
\end{center}

\paragraph{Initialized identifiers}
\label{sec:grust:id_usage:grsync:eq:initialized}

The judgment \inits{\Eq}{\IDs} collects the identifiers of signals that require memory. In \GRSync,
such stateful signals can only be initialized at the top level of a component or within a \kft{when}
equation, following rule \nameref{rule:grust:id_usage:grsync:eq:initialized:when}. The latter is
permitted because \kft{when} constructs are specifically designed for event-triggered state updates.
In contrast, \kft{match} equations describe control flow logic. Consequently, rule
\nameref{rule:grust:id_usage:grsync:eq:initialized:match} prohibits the initialization of new
memories within the branches of a \kft{match}.

\begin{center}
    \begin{minipage}{0.95\textwidth}
        \centering
        \begin{mathpar}
            %% Set of equations
            \inferINITEq
            { \left(\inits{\Eq}{\IDs}\right)^* }
            { \inits{\setEq}{\bigsqcup \IDs^*} }
            {Blck}
            \label{rule:grust:id_usage:grsync:eq:initialized:eqs}

            %% Declarations
            \inferINITEq
            { }
            { \inits{\declEq{\many{(\ID \: \kft{:} \: \STY)}}{\Expr}}{\emptyset} }
            {Decl}
            \label{rule:grust:id_usage:grsync:eq:initialized:decl}

            %% Definition
            \inferINITEq
            { }
            { \inits{\defEq{\many{\ID}}{\Expr}}{\emptyset} }
            {Def}
            \label{rule:grust:id_usage:grsync:eq:initialized:def}

            %% Initialization
            \inferINITEq
            { }
            { \inits{\initEq{\ID}{\CST}}{\{\lastExpr{\ID}\}} }
            {Init}
            \label{rule:grust:id_usage:grsync:eq:initialized:init}
        \end{mathpar}
    \end{minipage}

    \begin{minipage}{0.99\textwidth}
        \centering
        \begin{mathpar}
            %% Match
            \inferINITEq
            { \left(\inits{\setEq[k]}{\emptyset}\right)^* }
            { \inits{\matchX{\Expr}{\many{(\normBr{\gPat{\Pat_k}{\Expr_k}}{\setEq[k]})}}}{\emptyset} }
            {Mtch}
            \label{rule:grust:id_usage:grsync:eq:initialized:match}

            %% When
            \inferINITEq
            { \IDs = \{(\lastExpr{\ID})^*\} \\ \left(\inits{\setEq[k]}{\emptyset}\right)^* }
            { \inits{
                    \whenEq{\initBr{\many{(\ID \: \kfm{=} \: \CST\kft{;})}}}
                    {\many{(\normBr{\gPat{\EPat_k}{\Expr_k}}{\setEq[k]})}}
                }{\IDs} }
            {When}
            \label{rule:grust:id_usage:grsync:eq:initialized:when}
        \end{mathpar}
    \end{minipage}
\end{center}

\paragraph{Defined identifiers}
\label{sec:grust:id_usage:grsync:eq:defined}

The judgment \defines{\Eq}{\IDs} collects all identifiers assigned a value in an equation \Eq.

For a \kft{match} equation, rule \nameref{rule:grust:id_usage:grsync:eq:defined:match} requires the
set of defined identifiers to be identical across all branches. This aligns with the role of a
\kft{match}: it guarantees that a complete and consistent set of outputs is always produced,
regardless of which branch is executed.

In contrast, a \kft{when} equation enables selective updates. Rule
\nameref{rule:grust:id_usage:grsync:eq:defined:when} defines its set of identifiers as the union of
those in its branches. This mechanism allows a branch to specify only the flows it needs to change.
If a signal is not defined in the active branch, it retains its previous value, while an undefined
event is not emitted.

\begin{center}
    \begin{minipage}{0.95\textwidth}
        \centering
        \begin{mathpar}
            %% Set of equations
            \inferDEFEq
            { \left(\defines{\Eq}{\IDs}\right)^* }
            { \defines{\setEq}{\bigsqcup \IDs^*} }
            {Blck}
            \label{rule:grust:id_usage:grsync:eq:defined:eqs}

            %% Declarations
            \inferDEFEq
            { }
            { \defines{\declEq{\many{(\ID \: \kft{:} \: \STY)}}{\Expr}}{\{\ID^*\}} }
            {Decl}
            \label{rule:grust:id_usage:grsync:eq:defined:decl}

            %% Definition
            \inferDEFEq
            { }
            { \defines{\defEq{\many{\ID}}{\Expr}}{\{\ID^*\}} }
            {Def}
            \label{rule:grust:id_usage:grsync:eq:defined:def}

            %% Initialization
            \inferDEFEq
            { }
            { \defines{\initEq{\ID}{\CST}}{\emptyset} }
            {Init}
            \label{rule:grust:id_usage:grsync:eq:defined:init}
        \end{mathpar}
    \end{minipage}

    \begin{minipage}{0.95\textwidth}
        \centering
        \begin{mathpar}
            %% Match
            \inferDEFEq
            { \left(\defines{\setEq[k]}{\{\ID^*\}}\right)^* }
            { \defines{\matchX{\Expr}{\many{(\normBr{\gPat{\Pat_k}{\Expr_k}}{\setEq[k]})}}}{\{\ID^*\}} }
            {Mtch}
            \label{rule:grust:id_usage:grsync:eq:defined:match}

            %% When
            \inferDEFEq
            { \left(\defines{\setEq[k]}{\IDs}\right)^* }
            { \defines{
                    \mathfoot{\whenEq{\initBr{\many{(\ID \: \kfm{=} \: \CST\kft{;})}}}
                        {\many{(\normBr{\gPat{\EPat_k}{\Expr_k}}{\setEq[k]})}}}
                }{\bigcup \IDs^*} }
            {When}
            \label{rule:grust:id_usage:grsync:eq:defined:when}
        \end{mathpar}
    \end{minipage}
\end{center}

\subsubsection{Expressions}
\label{sec:grust:id_usage:grsync:expr}

The judgment \exprIdUsage{\G}{\Expr} verifies that all identifiers in an expression \Expr are used
correctly according to the context \G.

The simplest expressions involve constants, identifiers, and memory access. Rule
\nameref{rule:grust:id_usage:grsync:expr:const} states that a constant is always valid. Rule
\nameref{rule:grust:id_usage:grsync:expr:id} requires that an identifier \ID be present in the
context. Rule \nameref{rule:grust:id_usage:grsync:expr:last} allows accessing the previous value of
a signal with \lastExpr{\ID} only if the identifier has \memoryScope, ensuring it is a stateful
signal.

\begin{center}
    \begin{minipage}{0.9\textwidth}
        \centering
        \begin{mathpar}
            %% Constant
            \inferIDE
            {  }
            { \exprIdUsage{\G}{\CST} }
            {Cst}
            \label{rule:grust:id_usage:grsync:expr:const}

            %% Identifier
            \inferIDE
            { \G[\ID] = \scope }
            { \exprIdUsage{\G}{\ID} }
            {ID}
            \label{rule:grust:id_usage:grsync:expr:id}

            %% Last
            \inferIDE
            { \G[\lastExpr{\ID}] = \memoryScope }
            { \exprIdUsage{\G}{\lastExpr{\ID}} }
            {Last}
            \label{rule:grust:id_usage:grsync:expr:last}
        \end{mathpar}
    \end{minipage}
\end{center}

For composite expressions, validation is recursive. Rules
\nameref{rule:grust:id_usage:grsync:expr:emit} and \nameref{rule:grust:id_usage:grsync:expr:op}
state that an event emission or an $n$-ary operation is valid if all its sub-expressions are valid.

\begin{center}
    \begin{minipage}{0.9\textwidth}
        \centering
        \begin{mathpar}
            %% Emit
            \inferIDE
            { \exprIdUsage{\G}{\Expr} }
            { \exprIdUsage{\G}{\emitExpr{\Expr}} }
            {Emit}
            \label{rule:grust:id_usage:grsync:expr:emit}

            %% $n$-ary operation
            \inferIDE
            { \left(\exprIdUsage{\G}{\Expr}\right)^* }
            { \exprIdUsage{\G}{\opExpr{\many{\Expr}}} }
            {Op}
            \label{rule:grust:id_usage:grsync:expr:op}
        \end{mathpar}
    \end{minipage}
\end{center}

Function and component applications are handled by rules
\nameref{rule:grust:id_usage:grsync:expr:funapp} and
\nameref{rule:grust:id_usage:grsync:expr:compapp}. These rules require that the called function or
component is defined in the context and that all argument expressions are valid.

\begin{center}
    \begin{minipage}{0.99\textwidth}
        \centering
        \begin{mathpar}
            %% Application
            % Function
            \inferIDE
            {
                \G[\ID_f] = \funItem{\ID_f}{\many{(\ID_i \: \kft{:} \: \TY_i)}}{\TY_o}{\many{\Decl}} \\\\
                \left(\exprIdUsage{\G}{\Expr}\right)^*
            }
            { \exprIdUsage{\G}{\appExpr{\ID_f}{\many{\Expr}}} }
            {AppF}
            \label{rule:grust:id_usage:grsync:expr:funapp}

            % Component
            \inferIDE
            {
                \G[\ID_c] = \compItem{\ID_c}{\many{(\ID_i \: \kft{:} \: \STY_i)}}{\many{(\ID_o \: \kft{:} \: \STY_o)}}{\many{\Eq}} \\\\
                \left(\exprIdUsage{\G}{\Expr}\right)^*
            }
            { \exprIdUsage{\G}{\appExpr{\ID_c}{\many{\Expr}}} }
            {AppC}
            \label{rule:grust:id_usage:grsync:expr:compapp}
        \end{mathpar}
    \end{minipage}
\end{center}

For a \kft{match} expression, rule \nameref{rule:grust:id_usage:grsync:expr:match} imposes several
conditions. First, the matched expression \Expr must be valid in the current context \G. Then, for
each branch, the identifiers declared in the pattern $\Pat_k$ are added to the context, creating a
local context $\G_k$. Both the guard and the body expressions of that branch must be valid in this
extended local context.

\begin{center}
    \begin{minipage}{0.9\textwidth}
        \centering
        \begin{mathpar}
            %% Match
            \inferIDE
            {
                \exprIdUsage{\G}{\Expr} \\
                \left(\declares{\Pat_k}{\{\ID_k^*\}}\right)^* \\
                \left(\G_k = \G \ctxconcat [(\localScope / \ID_k)^*] \right) \\
                \left(\exprIdUsage{\G_k}{\Expr_k^g}\right)^* \\
                \left(\exprIdUsage{\G_k}{\Expr_k}\right)^*
            }
            { \exprIdUsage{\G}{\matchX{\Expr}{\many{(\normBr{\gPat{\Pat_k}{\Expr_k^g}}{\Expr_k})}}} }
            {Mtch}
            \label{rule:grust:id_usage:grsync:expr:match}
        \end{mathpar}
    \end{minipage}
\end{center}

\subsubsection{Patterns}
\label{sec:grust:id_usage:grsync:pat}

Patterns do not use identifiers from the context; instead, they declare new identifiers that are
local to the branch where they appear.

\paragraph{Declared identifiers}
\label{sec:grust:id_usage:grsync:pat:declared}

The judgment \declares{\Pat}{\IDs} collects the set of identifiers \IDs declared by a pattern
\Pat. The rules show that only identifier patterns contribute to this set. Constants, wildcards, and
other structural patterns do not declare any identifiers, while tuple patterns aggregate with a
disjoint union the identifiers from their sub-patterns.

\begin{center}
    \begin{minipage}{0.9\textwidth}
        \centering
        \begin{mathpar}
            %% Constant
            \inferDECLPat
            { }
            { \declares{\CST}{\emptyset} }
            {Cst}
            \label{rule:grust:id_usage:grsync:pat:declared:const}

            %% Identifier
            \inferDECLPat
            { }
            { \declares{\ID}{\{\ID\}} }
            {ID}
            \label{rule:grust:id_usage:grsync:pat:declared:id}

            %% Wildcard
            \inferDECLPat
            { }
            { \declares{\_}{\emptyset} }
            {Wild}
            \label{rule:grust:id_usage:grsync:pat:declared:wild}

            %% Tuple
            \inferDECLPat
            { \left(\declares{\Pat}{\IDs}\right)^* }
            { \declares{\tuple{\Pat}}{\bigsqcup \IDs^*} }
            {Tupl}
            \label{rule:grust:id_usage:grsync:pat:declared:tuple}
        \end{mathpar}
    \end{minipage}
\end{center}

\subsubsection{Event Patterns}
\label{sec:grust:id_usage:grsync:epat}

The judgment \epatIdUsage{\G}{\EPat} validates the use of identifiers within an event pattern \EPat.

Rule \nameref{rule:grust:id_usage:grsync:epat:event} handles event detection, ensuring the source
event $\ID_e$ is defined in the context. Rule \nameref{rule:grust:id_usage:grsync:epat:edge} handles
rising edges, requiring the underlying boolean expression \Expr to be valid. A tuple of event
patterns is valid if each constituent pattern is valid, as specified by rule
\nameref{rule:grust:id_usage:grsync:epat:tuple}.

\begin{center}
    \begin{minipage}{0.9\textwidth}
        \centering
        \begin{mathpar}
            %% Event detection
            \inferIDEPat
            { \G[\ID_e] = \scope }
            { \epatIdUsage{\G}{\ID \: \kfm{=} \: \opt{\ID_e}} }
            {Dtct}
            \label{rule:grust:id_usage:grsync:epat:event}

            %% Rising edge
            \inferIDEPat
            { \exprIdUsage{\G}{\Expr} }
            { \epatIdUsage{\G}{\Expr} }
            {Edge}
            \label{rule:grust:id_usage:grsync:epat:edge}

            %% Tuple
            \inferIDEPat
            { \left(\epatIdUsage{\G}{\EPat}\right)^* }
            { \epatIdUsage{\G}{\tuple{\EPat}} }
            {Tupl}
            \label{rule:grust:id_usage:grsync:epat:tuple}
        \end{mathpar}
    \end{minipage}
\end{center}

\paragraph{Declared identifiers}
\label{sec:grust:id_usage:grsync:epat:declared}

The judgment \declares{\EPat}{\IDs} collects identifiers declared by an event pattern \EPat. As
shown by the rules, only the event detection pattern declares a new identifier, which captures the
value carried by the event. Rising edge patterns do not declare identifiers, and tuple patterns
collect declarations from their sub-patterns using a disjoint union.

\begin{center}
    \begin{minipage}{0.9\textwidth}
        \centering
        \begin{mathpar}
            %% Event detection
            \inferDECLEPat
            { }
            { \declares{\ID \: \kfm{=} \: \opt{\ID_e}}{\{\ID\}} }
            {Dtct}
            \label{rule:grust:id_usage:grsync:epat:declared:event}

            %% Rising edge
            \inferDECLEPat
            { }
            { \declares{\Expr}{\emptyset} }
            {Edge}
            \label{rule:grust:id_usage:grsync:epat:declared:edge}

            %% Tuple
            \inferDECLEPat
            { \left(\declares{\EPat}{\IDs}\right)^* }
            { \declares{\tuple{\EPat}}{\bigsqcup \IDs^*} }
            {Tupl}
            \label{rule:grust:id_usage:grsync:epat:declared:tuple}
        \end{mathpar}
    \end{minipage}
\end{center}

\subsection{Summary}
This section has established the formal rules that govern identifier usage in \GRust. Using a set of
judgments, it has provided an unambiguous specification for how identifiers are scoped, declared,
and defined across the language. The logic handles the distinct scoping models of both \GRFRP
services, which use ordered lexical scoping, and \GRSync components, which employ a mutually
recursive structure. Key design choices, such as the requirement for \kft{match} branches to define
identical sets of identifiers, are formally captured by these rules to ensure consistency and
predictability.

This formal blueprint is essential, providing a sound foundation for the type system and the
semantic proofs presented in the following chapter. In the compiler prototype, described in
\Cref{sec:grust:implem:choices:architecture}, this specification is realized through a symbol table
that enforces these rules at compile time. While the prototype compiler validates these scoping
rules, it strategically delegates the check for pattern exhaustiveness to the underlying \Rust
compiler. Ensuring that \kft{match} statements are exhaustive is essential for the correctness of
the semantics of \GRust, and leveraging the robust analysis of the \Rust compiler is a sound
engineering choice. Because the structure of a \GRust \kft{match} maps almost directly to its \Rust
counterpart, any errors reported by the \Rust compiler remain clear and traceable. This approach
allows the \GRust compiler to focus on domain-specific checks, while guaranteeing that every
identifier is used correctly and without ambiguity.
